\begin{description}
\item[a)] Using the data from tables 5 to 9 and recalling that MK class II
  corresponds to giant stars while MK classes Ia and Iab correspond to
  supergiants, we easily obtain tables 10 to 12 and hence figures containing
  the plots of $E_{\mathrm{X-V}}/E_{\mathrm{B-V}}$ against $1/\lambda_X$ for
  both stars.

  \begin{center}
  Table 10 (10 points)

  \begin{tabular}{|l|c|c|c|c|c|c|c|c|c|c|}
  \hline
  Star & $\frac{B-V}{\mathrm{mag}}$ & $\frac{V-V}{\mathrm{mag}}$
    & $\frac{R-V}{\mathrm{mag}}$ & $\frac{I-V}{\mathrm{mag}}$
    & $\frac{J-V}{\mathrm{mag}}$ & $\frac{H-V}{\mathrm{mag}}$
    & $\frac{K-V}{\mathrm{mag}}$ & $\frac{L-V}{\mathrm{mag}}$
    & $\frac{M-V}{\mathrm{mag}}$ & $\frac{N-V}{\mathrm{mag}}$\\
  \hline
  HD 4817 & 1.9 & 0 & $-1.45$ & $-2.54$ & $-3.42$ & $-4.32$ & $-4.64$
    & $-4.86$ & $-4.59$ & --\\
  \hline
  HD 11092 & 2.09 & 0 & -- & -- & $-3.47$ & $-4.43$ & $-4.94$ & $-5.16$
    & $-4.92$ & $-5.13$\\
  \hline
  \end{tabular}
  \end{center}

  \begin{center}
  Table 11 (10 points)

  \begin{tabular}{|l|c|c|c|c|c|c|c|c|c|c|}
  \hline
  Star & $\frac{(B-V)_0}{\mathrm{mag}}$ & $\frac{(V-V)_0}{\mathrm{mag}}$
    & $\frac{(R-V)_0}{\mathrm{mag}}$ & $\frac{(I-V)_0}{\mathrm{mag}}$
    & $\frac{(J-V)_0}{\mathrm{mag}}$ & $\frac{(H-V)_0}{\mathrm{mag}}$
    & $\frac{(K-V)_0}{\mathrm{mag}}$ & $\frac{(L-V)_0}{\mathrm{mag}}$
    & $\frac{(M-V)_0}{\mathrm{mag}}$ & $\frac{(N-V)_0}{\mathrm{mag}}$\\
  \hline
  HD 4817 & 1.42 & 0 & $-1.13$ & $-1.96$ & $-2.41$ & $-3.14$ & $-3.25$
    & $-3.39$ & $-3.25$ & $-3.63$\\
  \hline
  HD 11092 & 1.42 & 0 & $-0.96$ & $-1.61$ & $-2.16$ & $-2.77$ & $-3.05$
    & $-3.22$ & $-3.08$ & $-3.02$\\
  \hline
  \end{tabular}
  \end{center}

  \begin{center}
  Table 12 (10 points)

  \begin{tabular}{|l|c|c|c|c|c|c|c|c|c|c|}
  \hline
  Star & $\frac{E_{\mathrm{B-V}}}{E_{\mathrm{B-V}}}$
    & $\frac{E_{\mathrm{V-V}}}{E_{\mathrm{B-V}}}$
    & $\frac{E_{\mathrm{R-V}}}{E_{\mathrm{B-V}}}$
    & $\frac{E_{\mathrm{I-V}}}{E_{\mathrm{B-V}}}$
    & $\frac{E_{\mathrm{J-V}}}{E_{\mathrm{B-V}}}$
    & $\frac{E_{\mathrm{H-V}}}{E_{\mathrm{B-V}}}$
    & $\frac{E_{\mathrm{K-V}}}{E_{\mathrm{B-V}}}$
    & $\frac{E_{\mathrm{L-V}}}{E_{\mathrm{B-V}}}$
    & $\frac{E_{\mathrm{M-V}}}{E_{\mathrm{B-V}}}$
    & $\frac{E_{\mathrm{N-V}}}{E_{\mathrm{B-V}}}$\\
  \hline
  HD 4817 & 1.00 & 0.00 & $-0.67$ & $-1.21$ & $-2.10$ & $-2.46$ & $-2.90$
    & $-3.06$ & $-2.79$ & --\\
  \hline
  HD 11092 & 1.00 & 0.00 & -- & -- & $-1.96$ & $-2.48$ & $-2.82$ & $-2.90$
    & $-2.75$ & $-3.15$\\
  \hline
  \end{tabular}
  \end{center}

  % FIGURE NEEDED: plot for HD 4817 (2015_da_sol.pdf p.8):
  % E(X-V)/E(B-V) against 1/lambda (1/nm), points rising from about -3 at
  % 1/lambda ~ 0.0003 to +1 at ~ 0.00225. (10 points)
  % FIGURE NEEDED: plot for HD 11092 (2015_da_sol.pdf p.8): same axes,
  % points rising from about -3.15 at 1/lambda ~ 0.0001 to +1 at ~ 0.00225.
  % (10 points)
\item[b)] We have that
  $E_{\lambda-\mathrm{V}} = A_\lambda - A_{\mathrm{V}}
  \rightarrow -A_{\mathrm{V}}$ as $\lambda\rightarrow\infty$, so the plot of
  $E_{\mathrm{X-V}}/E_{\mathrm{B-V}}$ against $1/\lambda_X$ intersects the
  vertical axis at
  \[ \frac{-A_{\mathrm{V}}}{E_{\mathrm{B-V}}} = -R_{\mathrm{V}}. \]
  \hfill(10 points)

  Hence $R_{\mathrm{V}}$ can be read off as minus the intersection of the
  plot with vertical axis. For the two stars we get Table 13 (the
  intersection points are obtained by fitting a curve to guide the eye,
  noting that as $1/\lambda\rightarrow0$, the curve becomes flat, as hinted
  at above).

  \begin{center}
  Table 13 (10 points)

  \begin{tabular}{|l|c|}
  \hline
  Star & $R_{\mathrm{V}}$\\
  \hline
  HD 4817 & 3.1\\
  \hline
  HD 11092 & 3.0\\
  \hline
  \end{tabular}
  \end{center}

  Next,
  \[ R_{\mathrm{r,r-i}} = \frac{A_{\mathrm{r}}}{E_{\mathrm{r-i}}}
     = \left(\frac{E_{\mathrm{r-V}}}{E_{\mathrm{B-V}}}
       + \frac{A_{\mathrm{V}}}{E_{\mathrm{B-V}}}\right)
       \frac{E_{\mathrm{B-V}}}{E_{\mathrm{r-i}}}
     = \frac{\dfrac{E_{\mathrm{r-V}}}{E_{\mathrm{B-V}}} + R_{\mathrm{V}}}
       {\dfrac{E_{\mathrm{r-V}}}{E_{\mathrm{B-V}}}
        - \dfrac{E_{\mathrm{i-V}}}{E_{\mathrm{B-V}}}}, \]
  where $E_{\mathrm{r-V}}/E_{\mathrm{B-V}}$ and
  $E_{\mathrm{i-V}}/E_{\mathrm{B-V}}$ can be read off from the plot (again,
  by fitting a curve to guide the eye). Hence we get Table 14.

  \begin{center}
  Table 14 (10 points)

  \begin{tabular}{|l|c|}
  \hline
  Star & $R_{\mathrm{r}}$\\
  \hline
  HD 4817 & 3.7\\
  \hline
  HD 11092 & 3.6\\
  \hline
  \end{tabular}
  \end{center}

  Hence we take the expected values to be $R_{\mathrm{V}} \approx 3.1$ and
  $R_{\mathrm{r}} \approx 3.7$. Note that the first value agrees with the
  widely used ratio of the total to selective extinction in filters B and V.
\item[c)] First let us find the apparent distance moduli $\mu_{\mathrm{r,i}}$
  in filters r and i. Reading off the fitted values e.g.\ at
  $\log(P/\mathrm{day}) = 1.6$ from figures 2 and 3 and substituting into
  the period-luminosity relations, we find $\mu_{\mathrm{r}} = 29.0$\,mag
  and $\mu_{\mathrm{i}} = 28.6$\,mag, so
  $E_{\mathrm{r-i}} = A_{\mathrm{r}} - A_{\mathrm{i}}
  = \mu_{\mathrm{r}} - \mu_{\mathrm{i}} = 0.4$\,mag and so
  $A_{\mathrm{r}} \approx 3.7E_{\mathrm{r-i}} = 1.5$\,mag. Hence the
  unreddened distance modulus is
  $\mu_0 = \mu_{\mathrm{r}} - A_{\mathrm{r}} = 27.5$\,mag and so we estimate
  the distance to IC 342 to be 3.2\,Mpc. \hfill(20 points)
\end{description}
